\documentclass[a4paper,11pt]{article}
\usepackage{exam-style-341}
\examinehead{341 农业综合知识三 · 网络技术分章练习}{第09章 网络安全}
\begin{document}

\unitcover{341 农业综合知识三 · 网络技术分章练习}{第09章\quad 计算机网络安全}{练习（共 100 分）}{约 60 分钟}{网络安全问题与攻击类型；对称密钥密码体制（DES／AES）；公开密钥密码体制（RSA）；数字签名；报文鉴别与散列函数；SSL／TLS 等安全协议；防火墙与入侵检测。}

\begin{knowmap}
\textbf{安全目标：}保密性（不泄露）· 完整性（不被篡改）· 可用性（不被阻断）；威胁分为被动攻击与主动攻击\\[2pt]
\textbf{被动攻击：}截获、流量分析，不改变数据，难以检测、重在预防；\textbf{主动攻击：}篡改、伪造、重放、拒绝服务，可检测、重在检测与恢复\\[2pt]
\textbf{对称密钥体制：}加密与解密使用同一密钥，速度快；典型算法 DES（56 位密钥）、AES；难点是密钥分配与管理\\[2pt]
\textbf{公钥密码体制：}公钥公开、私钥保密；公钥加密、私钥解密；RSA 基于大数分解困难；速度慢，常用于密钥交换与数字签名\\[2pt]
\textbf{鉴别与签名：}报文鉴别用散列函数（MD5 128 位、SHA-1 160 位）或 MAC；数字签名＝发送方私钥签名、接收方公钥验证，满足报文鉴别、完整性、不可否认\\[2pt]
\textbf{安全协议与系统安全：}IPsec（网络层）· SSL／TLS（运输层，HTTPS 443）· PGP（应用层）；防火墙（包过滤、应用网关）· 入侵检测（基于特征、基于异常）
\end{knowmap}

\begin{keyterms}
\textbf{被动攻击}：攻击者只窃听、分析通信内容而不改变数据，如截获、流量分析。\\
\textbf{主动攻击}：攻击者改变数据或伪造数据流，如篡改、伪造、重放、拒绝服务。\\
\textbf{对称密钥体制}：加密密钥与解密密钥相同的密码体制，典型算法有 DES、AES。\\
\textbf{DES}：数据加密标准，分组长度 64 位、有效密钥 56 位的对称密钥算法。\\
\textbf{AES}：高级加密标准，分组长度 128 位、密钥长度可变（128／192／256 位）的对称密钥算法。\\
\textbf{公钥密码体制}：加密密钥（公钥）公开、解密密钥（私钥）保密的密码体制，又称非对称密钥体制。\\
\textbf{RSA}：最著名的公钥密码算法，安全性基于大整数分解的困难性。\\
\textbf{数字签名}：发送方用私钥对报文（或其散列值）进行运算，接收方用发送方公钥验证，用于鉴别发送者与保证不可否认。\\
\textbf{报文鉴别}：验证报文的发送者身份和报文完整性，常用报文鉴别码 MAC 或数字签名实现。\\
\textbf{散列函数}：把任意长度报文映射为固定长度散列值的单向函数，如 MD5（128 位）、SHA-1（160 位）。\\
\textbf{SSL／TLS}：运输层安全协议，为应用层提供加密、鉴别与完整性保护，HTTPS 即 HTTP over TLS。\\
\textbf{防火墙}：部署在网络边界、按安全策略控制进出通信的访问控制设备，技术有包过滤和应用网关。\\
\textbf{入侵检测系统}：对网络流量进行深度检测、发现可疑入侵行为并报警的系统，分为基于特征和基于异常两类。\\
\textbf{拒绝服务攻击}：通过耗尽目标资源使其无法正常提供服务，破坏的是可用性。
\end{keyterms}

\begin{contrast}{对称密钥体制 / 公钥密码体制（每行依次为：密钥使用·典型算法·优缺点与用途）}
\citem{对称密钥体制}{加密与解密使用同一密钥}{DES、AES、IDEA}{速度快、适合大量数据加密；密钥分配与管理困难}
\citem{公钥密码体制}{公钥加密、私钥解密，或私钥签名、公钥验证}{RSA、ECC、ElGamal}{速度慢、适合密钥交换与数字签名；无需事先共享密钥}
\end{contrast}

\begin{contrast}{数字签名 / 报文鉴别（每行依次为：主要目的·使用的密钥·能否防止抵赖）}
\citem{数字签名}{鉴别发送者、保证报文完整、防止发送方抵赖}{发送方私钥签名，接收方用发送方公钥验证}{能防止抵赖}
\citem{报文鉴别}{验证发送者身份与报文完整性}{MAC 用双方共享的对称密钥}{不能防止发送方抵赖（双方都可能生成 MAC）}
\end{contrast}

\begin{contrast}{被动攻击 / 主动攻击（每行依次为：是否改变数据·典型手段·防范重点）}
\citem{被动攻击}{不改变数据，只窃听与分析}{截获、流量分析}{以防为主：加密、屏蔽、防泄漏}
\citem{主动攻击}{改变或伪造数据与数据流}{篡改、伪造、重放、拒绝服务}{以检测与恢复为主：鉴别、完整性校验、入侵检测、备份}
\end{contrast}

\qtypename{一、单项选择题}{10 题，每题 3 分，共 30 分}
\q 下列攻击中属于被动攻击的是（\quad）
\choicesline{A. 篡改报文}{B. 伪造报文}{C. 截获报文}{D. 拒绝服务}

\q DES 属于（\quad）
\choicesline{A. 对称密钥体制}{B. 公钥密码体制}{C. 散列函数}{D. 数字签名算法}

\q 在公钥密码体制中，A 要向 B 发送加密报文，应使用（\quad）加密。
\choicesline{A. A 的公钥}{B. A 的私钥}{C. B 的公钥}{D. B 的私钥}

\q 数字签名中，发送方签名使用（\quad），接收方验证签名使用（\quad）。
\choices{A. 发送方私钥，发送方公钥}{B. 发送方公钥，发送方私钥}
{C. 接收方私钥，接收方公钥}{D. 接收方公钥，接收方私钥}

\q MD5 散列函数输出的散列值长度是（\quad）
\choicesline{A. 64 位}{B. 128 位}{C. 160 位}{D. 256 位}

\q SSL／TLS 协议位于（\quad）
\choicesline{A. 网络层}{B. 运输层}{C. 会话层}{D. 应用层}

\q HTTPS 使用的默认端口是（\quad）
\choicesline{A. 80}{B. 443}{C. 8080}{D. 21}

\q 防火墙的主要功能是（\quad）
\choices{A. 检测并清除计算机病毒}{B. 对进出网络边界的通信实施访问控制}
{C. 对报文进行加密}{D. 为邮件提供加密存储}

\q 拒绝服务攻击主要破坏信息安全的（\quad）
\choicesline{A. 保密性}{B. 完整性}{C. 可用性}{D. 不可否认性}

\q RSA 算法的安全性主要基于（\quad）
\choices{A. 大整数分解的困难性}{B. 离散对数问题}{C. 对称密钥长度}{D. 散列函数的单向性}

\qtypename{二、多项选择题}{3 题，每题 4 分，共 12 分；每题有两个或两个以上正确项}
\q 下列属于主动攻击的有（\quad）
\choicesline{A. 篡改}{B. 伪造}{C. 重放}{D. 截获}

\q 下列属于对称密钥密码体制算法的是（\quad）
\choicesline{A. DES}{B. AES}{C. IDEA}{D. RSA}

\q 数字签名应满足的功能包括（\quad）
\choices{A. 报文鉴别}{B. 报文完整性}{C. 不可否认}{D. 提高传输速率}

\qtypename{三、判断题}{6 题，每题 2 分，共 12 分；正确的画“√”，错误的画“×”并写出改正后的表述}
\q 对称密钥体制中，加密密钥与解密密钥是相同的。\quad（\quad）\quad 改正：\fixblank

\q 公钥密码体制中，公钥可以公开，私钥必须保密。\quad（\quad）\quad 改正：\fixblank

\q 数字签名使用接收方的公钥对报文进行签名。\quad（\quad）\quad 改正：\fixblank

\q MD5 是一种对称加密算法。\quad（\quad）\quad 改正：\fixblank

\q 防火墙可以防止来自内部网络的攻击行为。\quad（\quad）\quad 改正：\fixblank

\q 报文鉴别可以同时验证报文的发送者身份和报文完整性。\quad（\quad）\quad 改正：\fixblank

\qtypename{四、计算题}{1 题，共 6 分；要求写出推导过程}
\q 在 RSA 密码体制中，设素数 $p=3$、$q=11$，取公开指数 $e=3$。请回答：
\begin{enumerate}[leftmargin=2em, itemsep=1pt, topsep=2pt]
  \item 求模数 $n$ 和欧拉函数 $\varphi(n)$。
  \item 求私钥 $d$，使 $e\cdot d\equiv 1\pmod{\varphi(n)}$。
  \item 对明文 $M=5$ 加密，求密文 $C$。
  \item 验证用私钥 $d$ 解密可以还原明文 $M$。
\end{enumerate}
\anslines{5cm}

\newpage
\qtypename{五、简答题}{2 题，每题 10 分，共 20 分}
\q 简述对称密钥体制与公开密钥体制的主要区别，并说明各自的典型算法与主要用途。\anslines{5cm}

\q 简述数字签名的基本原理及其应满足的三项功能，并说明实际中为什么先对报文做散列（Hash）再签名。\anslines{5cm}

\qtypename{六、综合题}{1 题，共 20 分}
\q 某农业企业建设“农产品质量安全溯源＋网上销售”系统，包含 Web 服务器、数据库、田间物联网采集设备和移动支付功能。请回答：
\begin{enumerate}[leftmargin=2em, itemsep=1pt, topsep=2pt]
  \item 该系统面临的主要安全威胁有哪些？请按被动攻击和主动攻击分类列举（6 分）。
  \item 为保护用户访问网站和在线支付的安全，应采用哪些安全协议与技术？简述其作用（8 分）。
  \item 为防止溯源数据被篡改，可设计怎样的技术方案？说明其原理（6 分）。
\end{enumerate}
\anslines{8cm}

\end{document}
